\documentclass[11pt]{article}
\usepackage{booktabs}
\usepackage{multirow}
\usepackage{tabularx}

\begin{document}

\section{Garden-path effects by construction}

Table~\ref{tab:npz} breaks the effect down by construction and condition.
Standard deviations are over three seeds. The NP/Z contrast is reliably larger
than NP/S in both models.

\begin{table*}[t]
\centering
\begin{tabularx}{\textwidth}{llcc}
\toprule
\multirow{2}{*}{Model} & \multirow{2}{*}{Condition} & \multicolumn{2}{c}{Surprisal $\Delta$} \\
\cmidrule(lr){3-4}
 & & NP/Z & NP/S \\
\midrule
\multirow{2}{*}{GPT-2-large} & Ambiguous & $1.42 \pm 0.21$ & $0.88 \pm 0.19$ \\
 & Unambiguous & $\textbf{0.31} \pm 0.11$ & $0.24 \pm 0.09$ \\
\midrule
\multirow{2}{*}{Pythia-1.4B} & Ambiguous & $1.19 \pm 0.24$ & $0.79 \pm 0.22$ \\
 & Unambiguous & $0.28 \pm 0.13$ & $0.21 \pm 0.10$ \\
\bottomrule
\end{tabularx}
\caption{Surprisal difference in bits at the disambiguating region, mean $\pm$
standard deviation over three seeds. Bold marks the smallest effect per column.}
\label{tab:npz}
\end{table*}

\end{document}
